微积分的每一句话都在讲``某个量随另一个量怎样变''。可是在动笔算变化率之前，得先把三件事钉死：谁依赖谁、允许喂进去的输入有哪些、输出被承诺落在什么范围里。这三件事含混，后面的极限、导数、积分全都悬在半空。

这听起来像形式手续，实际会直接改变判断。一份接口文档写着``模型输出一个 \(0\) 到 \(1\) 之间的分数''，它说的是输出\emph{可以}落在哪里；真跑一遍数据才发现分数集中在 \(0.02\) 与 \(0.87\) 之间，这才是输出\emph{实际}落在哪里。把两者混为一谈，就会去调一个永远触发不了的阈值。同样，\(\sqrt{1-x^2}\) 的定义域是 \([-1,1]\)，理由不是``开方不能取负''这条口诀，而是超出这一段之后输出根本不再是实数。定义域是规则的一部分，不是事后补上的免责声明。

本章沿着``一个函数可以被怎样观察''往下走。先把函数拆成定义域、陪域与对应规则，看清图像能告诉你什么、又会在哪里骗你；再看怎样把简单函数搭成复杂函数，以及这条流水线能不能倒着走回去；最后给出两族在后续每一章都会重新露面的基本变化模式——指数刻画固定倍数的增长，三角函数刻画绕一圈回到原处的重复。

\begin{itemize}
\tightlist
\item
  \hyperref[sec:02-Calculus/01-Functions-and-Precalculus/01]{``从实数轴到函数''}：把``接近''变成可以检查的量，并把陪域的承诺和值域的事实分开。
\item
  \hyperref[sec:02-Calculus/01-Functions-and-Precalculus/02]{``图像变换与函数复合''}：平移缩放不过是最简单的复合；顺序一换，结果和定义域一起变。
\item
  \hyperref[sec:02-Calculus/01-Functions-and-Precalculus/03]{``反函数与可逆性''}：什么时候能从输出还原输入，什么时候信息已经在中途被撞没了。
\item
  \hyperref[sec:02-Calculus/01-Functions-and-Precalculus/04]{``指数与对数：相同相对比例的变化''}：固定倍数的变化怎样写，对数坐标为什么能把它拉直。
\item
  \hyperref[sec:02-Calculus/01-Functions-and-Precalculus/05]{``三角函数与周期：用旋转描述重复变化''}：周期不是波形的性质，而是转满一圈回到同一点的后果。
\end{itemize}

五篇按依赖排列，顺读最省力；若从某个具体困惑进来，也可以直接翻最接近的一篇，再沿篇内链接回补。读每一篇时不妨都追问同样两句：输入允许取什么，输出实际能到哪里。这两句话在后面每一次讨论极限、连续与可导时，都会换一副面孔再出现。
